% =============================================================================
%  Notas de clase — Sesión 02: Introducción a C++
%  Programación Avanzada — Universidad Panamericana
%  Compilar: pdflatex sesion02_introduccion_cpp.tex
% =============================================================================
\documentclass[12pt, oneside]{article}
\usepackage{progavanzada}   % estilo "for dummies" alumno

\begin{document}

\portada{Sesión 02}{Introducción a C++}{2026}

\tableofcontents
\newpage

% =============================================================================
\section{¿Qué es C++?}
% =============================================================================

C++ es un lenguaje de programación de \textbf{propósito general} creado por
\textbf{Bjarne Stroustrup} a principios de los años 80 como extensión de C.
Su nombre alude al operador de incremento: \textit{C mejorado}.

\begin{recuerda}
  C++ es un lenguaje \textbf{compilado} y \textbf{tipado estático}: el código
  se traduce a lenguaje máquina antes de ejecutarse, y los tipos de las variables
  se verifican en tiempo de compilación.
\end{recuerda}

\subsection{Características principales}

\begin{itemize}
  \item \textbf{Compilado} — genera ejecutables rápidos.
  \item \textbf{Tipado estático} — errores detectados antes de ejecutar.
  \item \textbf{Acceso a memoria} — control directo mediante apuntadores.
  \item \textbf{Multi-paradigma} — soporta programación estructurada, OO y genérica.
  \item \textbf{Estándar vivo} — C++11, C++14, C++17, C++20, C++23.
\end{itemize}

% =============================================================================
\section{Dónde buscar documentación}
% =============================================================================

\begin{center}
\begin{tabular}{lp{9cm}}
  \toprule
  \textbf{Recurso} & \textbf{Cuándo usarlo} \\
  \midrule
  \href{https://en.cppreference.com}{cppreference.com}
    & Referencia técnica completa. La más precisa y actualizada. \\
  \href{https://cplusplus.com/reference}{cplusplus.com}
    & Más amigable para principiantes. \\
  \href{https://isocpp.org}{isocpp.org}
    & Sitio oficial del comité de estándares. \\
  \href{https://www.learncpp.com}{learncpp.com}
    & Tutorial paso a paso, excelente para aprender. \\
  \bottomrule
\end{tabular}
\end{center}

\begin{consejo}
  Ante cualquier duda sobre una función o clase de la biblioteca estándar,
  consulta primero \href{https://en.cppreference.com}{\textbf{cppreference.com}}.
  Aprender a leer documentación técnica es una habilidad esencial.
\end{consejo}

% =============================================================================
\section{Compiladores}
% =============================================================================

\subsection{g++ — GNU Compiler Collection}

El compilador de referencia en Linux y macOS.
Disponible en Windows con \href{https://www.mingw-w64.org/}{MinGW} o WSL.

\subsection{clang++ — LLVM}

Mensajes de error más claros y descriptivos. Preferido en macOS.

\subsection{cling — Intérprete interactivo}

No es un compilador tradicional: ejecuta C++ línea a línea.
Es el motor detrás de \textbf{xeus-cling}, el kernel C++ de Jupyter.

\begin{tecnico}
  \textbf{Compilar} (g++, clang++) convierte todo el programa a un ejecutable.
  \textbf{Interpretar} (cling) evalúa el código celda a celda, como Python.
  En producción siempre se usa un compilador real.
\end{tecnico}

% =============================================================================
\section{Entornos de trabajo}
% =============================================================================

\begin{itemize}
  \item \textbf{Consola / Terminal} — universal, solo necesitas el compilador.
  \item \textbf{Visual Studio Code} (\url{https://code.visualstudio.com}) —
        editor ligero con extensión C/C++ de Microsoft.
  \item \textbf{Xeus-cling en Jupyter} — interactivo, ideal para aprender.
  \item \textbf{Google Cloud Shell} (\url{https://shell.cloud.google.com}) —
        Linux completo en el navegador, \cpp{g++} preinstalado.
\end{itemize}

% =============================================================================
\section{Compilar con g++ desde la consola}
% =============================================================================

\begin{ejemplo}[Flujo de compilación básico]
\begin{codigo}[language=bash, numbers=none]
# 1. Verificar instalacion
g++ --version

# 2. Compilar
g++ -std=c++17 main.cpp -o hola

# 3. Ejecutar
./hola

# 4. Con avisos activados (buena practica)
g++ -std=c++17 -Wall -Wextra main.cpp -o hola
\end{codigo}
\end{ejemplo}

\begin{advertencia}
  Si el código tiene errores, g++ \textbf{no genera el ejecutable}. Lee el
  mensaje de error: indica el archivo, el número de línea y la descripción
  del problema. No ignores los avisos (\textit{warnings}).
\end{advertencia}

% =============================================================================
\section{Estructura básica de un programa C++}
% =============================================================================

\begin{codigo}[caption={main.cpp -- Hola Mundo}]
#include <iostream>
using namespace std;

int main(){
    cout << "Hello World";
}
\end{codigo}

\begin{center}
\begin{tabular}{lp{9.5cm}}
  \toprule
  \textbf{Línea} & \textbf{Significado} \\
  \midrule
  \cpp{\#include <iostream>}
    & Incluye la biblioteca de entrada/salida. Sin esta, \cpp{cout} no existe. \\
  \cpp{using namespace std;}
    & Permite usar \cpp{cout}, \cpp{cin}, \cpp{endl} sin prefijo \cpp{std::}. \\
  \cpp{int main()\{ \}}
    & Punto de entrada: todo programa C++ debe tener exactamente una función \cpp{main}. \\
  \cpp{cout << "...";}
    & Imprime texto en la consola. \cpp{<<} es el operador de inserción de flujo. \\
  \bottomrule
\end{tabular}
\end{center}

% =============================================================================
\section{Entrada y salida: cin y cout}
% =============================================================================

\begin{codigo}[caption={main3.cpp -- Saludo personalizado}]
#include <iostream>
#include <string>
using namespace std;

int main(){
    string nombre;
    cout << "Introduce tu nombre: ";
    cin >> nombre;
    cout << "Hello " << nombre << endl;
}
\end{codigo}

\begin{consejo}
  \cpp{cin >>} lee \textbf{una palabra} (hasta el primer espacio). Para leer
  una línea completa usa \cpp{getline(cin, nombre)}.
\end{consejo}

\begin{advertencia}
  En el cuaderno Jupyter (xeus-cling), \cpp{cin} \textbf{no funciona de forma
  interactiva}. Asigna el valor directamente a la variable para demostrar
  el mismo concepto dentro del cuaderno.
\end{advertencia}

% =============================================================================
\section{Jueces en línea}
% =============================================================================

Un \textbf{juez en línea} (\textit{online judge}) es una plataforma web que
recibe tu código, lo compila en sus servidores, lo ejecuta contra casos de
prueba ocultos y te dice si tu solución es correcta — sin que nadie tenga que
revisarla a mano. Es la forma estándar de practicar y competir en programación
a nivel mundial.

\subsection{¿Cómo funciona?}

\begin{enumerate}
  \item Lees el enunciado del problema (entrada esperada, salida esperada).
  \item Escribes tu solución en C++ (o el lenguaje permitido).
  \item Subes el código al juez.
  \item El juez lo compila y lo corre contra decenas de casos de prueba.
  \item Recibes un veredicto:
\end{enumerate}

\begin{center}
\begin{tabular}{ll}
  \toprule
  \textbf{Veredicto} & \textbf{Significado} \\
  \midrule
  \texttt{AC} — Accepted           & Tu solución es correcta en todos los casos \\
  \texttt{WA} — Wrong Answer       & Falla en al menos un caso de prueba \\
  \texttt{TLE} — Time Limit Exceeded & Correcta pero demasiado lenta \\
  \texttt{MLE} — Memory Limit Exceeded & Usa demasiada memoria \\
  \texttt{CE} — Compilation Error  & Tu código no compila \\
  \texttt{RE} — Runtime Error      & El programa crashea durante la ejecución \\
  \bottomrule
\end{tabular}
\end{center}

\begin{recuerda}
  El veredicto \texttt{AC} no significa que tu código sea perfecto — significa
  que pasó los casos de prueba del juez. Un buen programador se pregunta además:
  ¿es legible? ¿es eficiente? ¿qué casos extremos no consideré?
\end{recuerda}

\subsection{omegaUp — El juez mexicano}

\begin{center}
  \url{https://omegaup.com}
\end{center}

omegaUp es una plataforma de programación competitiva creada en \textbf{México},
con sede en la Ciudad de México y respaldada por egresados de instituciones
mexicanas. Es el juez oficial de la \textbf{Olimpiada Mexicana de Informática
(OMI)} y cuenta con miles de problemas en español, incluyendo material
pedagógico pensado para estudiantes latinoamericanos.

\begin{tecnico}
  omegaUp fue fundado en 2011 por Luis Héctor Chávez y un equipo de
  programadores mexicanos con el objetivo de democratizar el acceso a la
  programación competitiva en América Latina. Hoy es una organización sin
  fines de lucro con más de 100\,000 usuarios registrados y es utilizado por
  preparatorias y universidades en toda la región.\\
  \textit{Referencia:} \url{https://omegaup.com/about}
\end{tecnico}

\subsection{VJudge — Agregador de jueces internacionales}

\begin{center}
  \url{https://vjudge.net}
\end{center}

VJudge no es un juez propio: es un \textbf{agregador} que reúne problemas de
los jueces más importantes del mundo (Codeforces, UVa, POJ, SPOJ, LeetCode,
entre otros) en una sola interfaz. Permite crear \textit{contests} privados
con problemas de múltiples fuentes, ideal para práctica en clase.

\begin{center}
\begin{tabular}{ll}
  \toprule
  \textbf{Plataforma} & \textbf{Enfoque} \\
  \midrule
  omegaUp   & Español, material pedagógico, OMI, contexto mexicano \\
  VJudge    & Acceso a miles de problemas internacionales desde una cuenta \\
  Codeforces & Competencias semanales, ranking mundial \\
  LeetCode  & Entrevistas técnicas, empresas tecnológicas \\
  \bottomrule
\end{tabular}
\end{center}

\subsection{Tarea: crea tus cuentas}

Antes de la siguiente sesión, registra una cuenta en ambas plataformas.
Usa un nombre de usuario que puedas compartir en clase (el profesor lo
usará para enviar problemas).

\begin{center}
\begin{tabular}{lll}
  \toprule
  \textbf{Plataforma} & \textbf{URL de registro} & \textbf{Mi usuario} \\
  \midrule
  omegaUp & \url{https://omegaup.com} & \hspace{3cm} \\[8pt]
  VJudge  & \url{https://vjudge.net}  & \hspace{3cm} \\
  \bottomrule
\end{tabular}
\end{center}

\begin{consejo}
  Usa el \textbf{mismo nombre de usuario} en ambas plataformas si está
  disponible. Facilita que el profesor te identifique y que construyas
  un perfil reconocible en la comunidad de programación competitiva.
\end{consejo}

% =============================================================================
\section{Resumen}
% =============================================================================

\begin{recuerda}
\begin{itemize}
  \item C++ es compilado, tipado estático, con acceso directo a memoria.
  \item Documentación principal: \href{https://en.cppreference.com}{cppreference.com}.
  \item Compilar desde consola: \cpp{g++ -std=c++17 archivo.cpp -o ejecutable}
  \item Todo programa necesita una función \cpp{main()}; en xeus-cling no.
  \item \cpp{cout <<} imprime; \cpp{cin >>} lee (solo en consola).
\end{itemize}
\end{recuerda}

\end{document}
